\subsection{라이브 스크리닝 파일럿 (pilot-only)}
별도 스크리닝 파일럿은 Codex CLI를 통해 gpt-5.6-sol을 호출했다. 셀마다 seed로 구분한 제안 호출 5회와 수리 예산 $K=1$을 사용했고, 각 호출의 최초 후보 하나를 블라인드 동일 후보 retry와 $\rho$가 공유했다. Hosted revision은 고정되지 않았고 token 회계는 제공되지 않았으며 seed는 hosted sampler를 제어하지 않으므로 이 호출들은 독립 무작위 표본이 아니라 준반복이다. 따라서 추론 통계나 모델 순위를 보고하지 않는다.

표~\ref{tab:live-screening}에서 유도 우위는 의도적으로 guided-repairable 오류가 나오도록 구성한 \texttt{signal-repair-v2}의 policy-blind 셀에서만 나타났다. 최초 후보 5/5가 모두 무효였고 $\rho$는 5/5, 블라인드는 0/5를 commit했다. 나머지 현재 셀 4개는 유도 우위를 보이지 않았다. 세 셀은 최초 후보가 이미 유효했고, 동결 기저의 goal-directed 셀은 수리 불가 quest-stage regression이었다. 현재 셀의 모든 noncommit arm 결과 15/15가 prior-state hash를 보존했다. 이 결과는 해당 오류 regime에서의 mechanism transfer만 지지하며 모집단 효능, 모델 승격, 표본 효율 일반화는 지지하지 않는다. 따라서 \texttt{C-PILOT-007}과 \texttt{C-PILOT-008}은 pilot-only이고 \texttt{C-RESULT-003}은 \texttt{TODO-RESULT}로 남는다. 폐기된 v1 진단에서는 동반 오류가 최대 3개에서 수리 불가 knowledge 오류 1개로 줄었지만 commit되지 않아, 오류 하나라도 수리 불가이면 부분 수리가 합성되지 않음을 보여 주었다.
